\section{Synthesis: Design Principles for Geographic Representation Systems}
\label{sec:synthesis}

The five lessons from the E-track experiments support a set of design principles
that apply beyond the specific systems studied.  This section synthesizes those
principles and derives their implications for the constraint hierarchy that
governs single-member district drawing.

\subsection{The Three Design Tensions}

The E-track experiments make visible three fundamental tensions in geographic
representation system design that are present in all systems but visible only
when systems are compared systematically:

\subparagraph{Tension 1: Compactness versus Proportionality.}
The most pervasive tension is between geographic compactness and partisan
proportionality.  Lesson~5 quantifies this tension at the Pareto frontier:
0.015 PP per pp of proportionality gained.  Lesson~1 explains its cause:
geographic sorting makes Democrats geographically inefficient in single-member
winner-take-all systems.  Lesson~2 establishes that multi-member districts
partially but not fully relieve this tension by averaging over more geographic
heterogeneity.

This tension is not merely a feature of the US context.  \citet{rodden2019why}
documents analogous sorting patterns in the UK, Canada, and Australia.
Paper E.6 finds the same proportionality--compactness exchange in all four
Westminster-system countries studied, confirming that the tension is a property
of the geographic sorting--winner-take-all combination, not of US partisan
geography specifically.

\subparagraph{Tension 2: Community Preservation versus Compactness.}
Lesson~3 and Lesson~4 both surface a second tension: the desire to preserve
pre-existing political communities (counties, states) conflicts with geometric
compactness optimization.  State boundaries cost 12\% PP (Lesson~4);
county preservation at the 85\% level costs 18\% PP (from the E.2 results
\citep{dellaLibera2026e2}).  These costs are independent: preserving both
counties and states costs approximately 28\% PP relative to unconstrained
national optimization.

The community-preservation tension is normatively distinct from the
compactness--proportionality tension.  Proportionality is contested as a
redistricting goal (courts have declined to constitutionalize it); community
preservation is widely accepted as a legitimate goal (most state redistricting
statutes list it as a criterion).  This means that community-preservation
costs are costs that the political system has implicitly agreed to pay, whereas
compactness--proportionality trade-offs are costs that are being debated.

\subparagraph{Tension 3: Structural Design versus Drawing Optimization.}
The deepest tension is not between two objectives on the Pareto frontier but
between two different levels of reform.  Lessons~1 and 2 together establish
that within any single-member geographic system, drawing optimization
(algorithmic redistricting, independent commissions, court supervision) can
eliminate process-level distortions but cannot eliminate structure-level
distortions arising from geographic sorting.  Structural reform --- multi-member
districts --- is required to materially change the seat--vote relationship.

This tension implies that the redistricting reform debate has been conducted
partly at the wrong level.  The bulk of litigation, legislative effort, and
research attention has focused on drawing-process reform.  The structural reform
that would actually change the seat--vote relationship (multi-member districts)
has received far less policy attention and far less empirical evaluation.
The E-track experiments are intended partly to correct this imbalance by
providing rigorous comparative evidence on what structural reform achieves.

\subsection{The Optimal Constraint Hierarchy for Single-Member Systems}

Given these three tensions, what is the optimal constraint hierarchy for
single-member district drawing?  We argue that the current practice ---
population equality first, contiguity second, compactness third, with
community preservation as a secondary criterion --- is optimal for the following
reasons.

\emph{Population equality first} is constitutionally mandated (\emph{Wesberry}
\citep{wesberry1964}) and electorally foundational.  No other criterion can
justify systematic population inequality.

\emph{Contiguity second} follows from population equality: a non-contiguous
district is geometrically incoherent as a representation unit and creates
representational obligations that the representative cannot fulfill for
constituents with no geographic connection to the district's main body.

\emph{Compactness third} follows from contiguity: given population equality
and contiguity, the shape of the district is the primary remaining degree of
freedom.  Compactness optimization within this constraint yields the districts
most resistant to partisan manipulation \citep{polsby1991third} and most
geographically coherent as communities of representation.

\emph{Community preservation as secondary criterion} reflects the lesson that
community preservation is valuable but costs compactness.  Treating it as a
hard constraint (preserve all counties) is too expensive; treating it as a
soft criterion (preserve counties where doing so does not substantially reduce
compactness) is the right balance.

The key design lesson for single-member systems is that adding proportionality
as a criterion at any level of the hierarchy is inappropriate: it substitutes
a structural problem (geographic sorting) for a drawing problem, and attempts
to solve the structural problem through drawing will produce maps that are less
compact than what the geography permits without achieving the proportionality
goal reliably.

\subsection{The Case for Multi-Member Districts as Structural Reform}

The synthesis of the five lessons points consistently toward multi-member
districts as the only mechanism within a geographic, single-election
representation system that can materially improve both proportionality and
minority representation without sacrificing the constitutional advantages of
district-based representation (geographic coherence, constituent-representative
accountability, local focus).

The advantages of multi-member districts over single-member alternatives are:
\begin{enumerate}
  \item \textbf{Partial sorting escape}: 40\% gap reduction (three-member) to
        65\% gap reduction (five-member), compared to 0\% for any single-member
        drawing method (Lessons~1 and 2).
  \item \textbf{Minority representation improvement}: +18\% majority-minority
        seats without requiring geographic concentration (Lesson~2).
  \item \textbf{Constitutional availability}: Requires only federal statute
        to repeal the 2~U.S.C.~\S2c single-member mandate (E.0
        constitutional taxonomy).
  \item \textbf{Compactness compatibility}: Three-member districts achieve
        most of the proportionality gain at a compactness cost (22\% PP
        reduction) that is large but not catastrophic --- far smaller than the
        34\%--56\% cost of five-member or pure-proportionality systems.
\end{enumerate}

The disadvantages --- ballot complexity, weaker local ties, incumbent advantage
in large districts --- are real but are normative objections rather than
empirical disqualifiers.  The ballot-complexity objection is partially addressed
by ranked-choice voting experience in STV systems (Ireland, Australia), where
voter error rates are modest and declining \citep{farrell2011electoral}.

We do not advocate for multi-member districts as a policy position.  We observe
that the E-track research program, taken as a whole, provides a stronger
empirical foundation for the structural case for multi-member districts than
previously existed, and that this evidence should inform reform debates even
if normative judgments about ballot complexity and local accountability ultimately
determine the political outcome.
